% =====================================================================
% BMED365 (Spring 2027) — Lab 4: Kiwi Fruit 3D Seed Segmentation Challenge
% Printable README PDF (generated faithfully from README.md)
% Compile with: pdflatex Lab4-Kiwi-challenge-README.tex (run twice)
% =====================================================================
\documentclass[10pt,a4paper]{article}

\usepackage[utf8]{inputenc}
\usepackage[T1]{fontenc}
\usepackage[english]{babel}
\usepackage{lmodern}
\usepackage{microtype}
\usepackage{graphicx}
\usepackage{booktabs}
\usepackage{xltabular}
\usepackage{array}
\usepackage[table]{xcolor}
\usepackage{amsmath,amssymb}
\usepackage{enumitem}
\usepackage[margin=1.7cm,top=1.75cm,bottom=1.85cm,headheight=14pt]{geometry}
\usepackage{fancyhdr}
\usepackage{lastpage}
\usepackage{titlesec}
\usepackage[most]{tcolorbox}
\usepackage{float}
\usepackage[hyphens]{url}
\usepackage{hyperref}
\def\UrlBreaks{\do\/\do-\do_\do.}

% =====================================================================
% COLORS & HYPERREF (UiB / BMED365 House Style)
% =====================================================================
\definecolor{uibblue}{RGB}{0, 61, 115}
\definecolor{uibred}{RGB}{175, 28, 44}
\definecolor{kiwigreen}{RGB}{27, 138, 90}
\definecolor{amberwarn}{RGB}{195, 110, 0}
\definecolor{lightgray}{RGB}{246, 248, 250}
\definecolor{rowalt}{RGB}{243, 247, 252}
\definecolor{codebg}{RGB}{242, 244, 247}

\hypersetup{
    colorlinks=true,
    linkcolor=uibblue,
    citecolor=uibblue,
    urlcolor=uibred,
    pdftitle={Lab 4: The Kiwi Fruit 3D Seed Segmentation Challenge (BMED365 Spring 2027)},
    pdfauthor={Arvid Lundervold (UiB / MMIV) -- BMED365 Spring 2027}
}

% =====================================================================
% TYPOGRAPHY, HEADERS & CALLOUT BOXES
% =====================================================================
\setlength{\parindent}{0pt}
\setlength{\parskip}{0.42em}
\setlist[itemize]{topsep=2pt,itemsep=2pt,leftmargin=1.4em}
\setlist[enumerate]{topsep=2pt,itemsep=2.5pt,leftmargin=1.6em}

\newcolumntype{L}[1]{>{\raggedright\arraybackslash}p{#1}}
\newcolumntype{C}[1]{>{\centering\arraybackslash}p{#1}}
\newcolumntype{Y}{>{\raggedright\arraybackslash}X}

\titleformat{\section}{\Large\bfseries\color{uibblue}\raggedright}{\thesection}{0.6em}{}[\vspace{1pt}\color{uibblue!40}\titlerule]
\titleformat{\subsection}{\large\bfseries\color{uibblue!88!black}\raggedright}{\thesubsection}{0.55em}{}
\titlespacing*{\section}{0pt}{1.05em}{0.42em}
\titlespacing*{\subsection}{0pt}{0.75em}{0.28em}

\pagestyle{fancy}
\fancyhf{}
\fancyhead[L]{\footnotesize\textcolor{gray}{BMED365 (Spring 2027) $\cdot$ Lab 4: Computational Imaging}}
\fancyhead[R]{\footnotesize\textcolor{gray}{The Kiwi Fruit 3D Seed Segmentation Challenge}}
\fancyfoot[L]{\footnotesize\textcolor{gray}{University of Bergen \& MMIV $\cdot$ \href{https://github.com/arvidl/BMED365-2027/tree/main/Lab4-Comp-Imag/Kiwi-challenge}{\textcolor{gray}{github.com/arvidl/BMED365-2027}}}}
\fancyfoot[R]{\footnotesize\textcolor{gray}{Page \thepage\ of \pageref{LastPage}}}
\renewcommand{\headrulewidth}{0.4pt}
\renewcommand{\footrulewidth}{0.2pt}

\fancypagestyle{firstpage}{
    \fancyhf{}
    \fancyfoot[L]{\footnotesize\textcolor{gray}{University of Bergen \& MMIV $\cdot$ \href{https://github.com/arvidl/BMED365-2027/tree/main/Lab4-Comp-Imag/Kiwi-challenge}{\textcolor{gray}{github.com/arvidl/BMED365-2027}}}}
    \fancyfoot[R]{\footnotesize\textcolor{gray}{Page \thepage\ of \pageref{LastPage}}}
    \renewcommand{\headrulewidth}{0pt}
    \renewcommand{\footrulewidth}{0.2pt}
}

\newtcolorbox{infobox}[1][]{
    enhanced, breakable,
    colback=lightgray, colframe=uibblue,
    fonttitle=\bfseries\small, title={#1},
    boxrule=0.6pt, arc=2.5pt, left=7pt, right=7pt, top=5pt, bottom=5pt
}

\newtcolorbox{kiwibox}[1][]{
    enhanced, breakable,
    colback=kiwigreen!5, colframe=kiwigreen!85!black,
    fonttitle=\bfseries\small, title={#1},
    boxrule=0.6pt, arc=2.5pt, left=7pt, right=7pt, top=5pt, bottom=5pt
}

\newtcolorbox{protocolbox}[1][]{
    enhanced, breakable,
    colback=amberwarn!5, colframe=amberwarn!85!black,
    fonttitle=\bfseries\small, title={#1},
    boxrule=0.6pt, arc=2.5pt, left=7pt, right=7pt, top=5pt, bottom=5pt
}

\newcommand{\colablink}[1]{%
  \href{#1}{\tcbox[on line, arc=2pt, boxsep=0pt, left=4.5pt, right=4.5pt, top=2pt, bottom=2pt,
    colback=amberwarn!15, colframe=amberwarn!80!black, boxrule=0.45pt]{%
    \sffamily\bfseries\scriptsize\color{amberwarn!55!black}Colab}}%
}

\newcommand{\ghblob}{https://github.com/arvidl/BMED365-2027/blob/main}
\newcommand{\ghtree}{https://github.com/arvidl/BMED365-2027/tree/main}
\newcommand{\colabbase}{https://colab.research.google.com/github/arvidl/BMED365-2027/blob/main}
\newcommand{\imgbase}{.}

% =====================================================================
\begin{document}
\thispagestyle{firstpage}

% --- TITLE BANNER ---
\begin{tcolorbox}[
    enhanced, colback=uibblue, colframe=uibblue,
    arc=3.5pt, left=12pt, right=12pt, top=8pt, bottom=8pt
]
{\Large\bfseries\color{white} Lab 4: The Kiwi Fruit 3D Seed Segmentation Challenge}\\[3pt]
{\normalsize\color{white!90} BMED365 --- Computational Imaging, Modeling and AI in Biomedicine (Spring 2027)}\\[3pt]
{\footnotesize\color{white!80}
Arvid Lundervold (University of Bergen \& MMIV) $\;\cdot\;$
Runs locally with \texttt{uv} (\texttt{bmed365-2027} kernel) or in Google Colab $\;\cdot\;$
\href{\ghtree/Lab4-Comp-Imag/Kiwi-challenge}{\textcolor{white}{\underline{Repository Folder}}}
}
\end{tcolorbox}

\begin{center}
\vspace{2pt}
\includegraphics[width=0.60\textwidth]{\imgbase/assets/photo_mri_kiwi_20050219.jpg}\\[2pt]
{\footnotesize\textit{Optical cross-section vs.\ preclinical 7\,T $T_2$-weighted RARE MRI of an intact kiwifruit (17 February 2005, MIC / Vivarium, UiB).}}
\vspace{-2pt}
\end{center}

% =====================================================================
\section{Overview: One Simple Fruit, Six Weeks of Computational Thinking}
% =====================================================================

Everyone knows what the inside of a kiwifruit looks like. Yet teaching a computer to accurately \textbf{find, separate, and count every single seed in 3D} from a real Magnetic Resonance Imaging (MRI) scan turns out to be surprisingly challenging --- and closely mirrors real clinical tasks in biomedical imaging.

In \textbf{BMED365}, the Kiwi Challenge serves as a \textbf{recurring benchmark problem} (``red thread'') from the first lecture in January to \textbf{Lab 4 (Computational Imaging)} in February:

\begin{itemize}
    \item \textbf{Intuitive biology, real MRI physics:} No medical jargon is needed on Day~1, yet the 3D scan exhibits partial-volume blurring, intensity shading, closely packed instances, and dark fibrous structures that mimic seeds.
    \item \textbf{Clinical analogy:} Replace the kiwifruit with the \textbf{human brain} and the seeds with \textbf{multiple sclerosis (MS) white-matter plaques}, \textbf{cerebral microbleeds}, or \textbf{micrometastases} (or replace the MRI with 3D microscopy of \textbf{cell nuclei}). The computational questions are identical.
    \item \textbf{Methodological playground:} You can attack the exact same 3D dataset with rule-based filters, classical machine learning, graph/spatial priors, deep learning, vision foundation models, neurosymbolic rules, and AI coding agents.
\end{itemize}

\textbf{Slides \& visual teasers:}
\begin{itemize}
    \item \href{\ghblob/Lab4-Comp-Imag/Kiwi-challenge/Beamer/main.pdf}{\textbf{Lecture 1 Teaser Slides (Beamer PDF, 10 min)}} \textit{(source: \href{\ghblob/Lab4-Comp-Imag/Kiwi-challenge/Beamer/main.tex}{\texttt{./Beamer/main.tex}})}
    \item \href{https://docs.google.com/presentation/d/1B_E-Pcwfvj3hgG2o02vsxmez_Bm6x72uI5tzIu0dymU/edit?usp=sharing}{\textbf{Original Problem Formulation (Google Slides)}}
    \item \href{\ghblob/Lab4-Comp-Imag/Kiwi-challenge/assets/kiwi_7T_flythrough.gif}{\textbf{Animated 7T MRI Axial Flythrough (\texttt{kiwi\_7T\_flythrough.gif})}} \textit{(static mid-equatorial slice shown below)}
\end{itemize}

\begin{center}
\vspace{2pt}
\begin{minipage}[c]{0.20\textwidth}
    \centering
    \includegraphics[width=\linewidth]{\imgbase/data/kiwi_064.png}\\
    {\scriptsize\textit{Mid-axial slice 64 (\href{\ghblob/Lab4-Comp-Imag/Kiwi-challenge/assets/kiwi_7T_flythrough.gif}{animated GIF})}}
\end{minipage}
\hfill
\begin{minipage}[c]{0.77\textwidth}
    \centering
    \includegraphics[width=\linewidth]{\imgbase/assets/kiwi_lecture1_teaser_card.png}\\
    {\scriptsize\textit{Lecture 1 teaser card: 7\,T MRI slice, naive thresholding failure, and curated 3D Micro-GT sector.}}
\end{minipage}
\end{center}

% =====================================================================
\clearpage
\section{The Challenge Tasks}
% =====================================================================

Starting from the 3D 7\,T MRI scan of an intact kiwifruit:

\begin{enumerate}
    \item \textbf{3D Seed Segmentation \& Counting:} Segment the dark seeds from the surrounding fruit tissue and estimate the \textbf{total number of seeds} ($N$) in the kiwifruit. If your method has tunable parameters or probabilistic outputs, provide an \textbf{uncertainty / sensitivity range} for your count.
    \item \textbf{Biomedical Biomarker Analogy (``Lesion Load''):} Treat the whole kiwifruit as a ``brain'' and the segmented seeds as focal pathological lesions (e.g., \href{https://pubmed.ncbi.nlm.nih.gov/23084503/}{MS plaques}). Using the physical voxel dimensions of the scan, estimate:
    \begin{itemize}
        \item Whole-fruit volume $V_{\text{kiwi}}$ ($\text{mm}^3$ or $\text{cm}^3$)
        \item Total seed volume $V_{\text{seeds}}$ ($\text{mm}^3$) and mean single-seed volume ($\text{mm}^3$)
        \item \textbf{``Lesion load''} (\href{https://journals.plos.org/plosone/article?id=10.1371/journal.pone.0120754}{volume fraction}): $V_{\text{seeds}} / V_{\text{kiwi}}$
    \end{itemize}
    \item \textbf{Verification Without Complete Ground Truth:} Nobody manually segmented all $\sim$700--1,200 seeds across all 128 slices, and the physical fruit from 2005 was never dissected to count the seeds. \textbf{How can you verify whether your count and segmentation are trustworthy?} (See Section~\ref{sec:eval}, \hyperref[sec:eval]{Evaluation \& Semi-Verifiability}, below.)
\end{enumerate}

% =====================================================================
\section{Data \& Notebooks}
% =====================================================================

All data files are included directly in \href{\ghtree/Lab4-Comp-Imag/Kiwi-challenge/data}{\texttt{./data/}} --- no external downloads are required.

\subsection{Dataset Summary}

The 3D dataset was acquired on \textbf{17 February 2005} using the preclinical \textbf{7\,T \href{https://www.bruker.com/en/products-and-solutions/preclinical-imaging/mri/pharmascan-new.html}{Bruker Pharmascan}} (\texttt{Pharmascan300}) MRI scanner at the \textit{Molecular Imaging Center} (MIC / Vivarium), University of Bergen.

{\footnotesize
\renewcommand{\arraystretch}{1.28}
\rowcolors{2}{rowalt}{white}
\begin{xltabular}{\textwidth}{L{2.55cm} L{5.35cm} Y}
\toprule
\rowcolor{uibblue!12}\textbf{Property} & \textbf{Value} & \textbf{Why it matters for segmentation} \\
\midrule
\endfirsthead
\toprule
\rowcolor{uibblue!12}\textbf{Property} & \textbf{Value} & \textbf{Why it matters for segmentation} \\
\midrule
\endhead
\bottomrule
\endfoot
\textbf{Pulse sequence} &
2D RARE (Rapid Acquisition with Relaxation Enhancement) &
$T_2$-weighted contrast: water-rich locular gel is bright; seeds, dense central core (\textit{columella}), radial septa, and peel appear dark. \\
\textbf{Sequence parameters} &
\texttt{TR = 18,520.8\,ms}, \texttt{TE = 63.9\,ms}, \texttt{ETL = 8}, \texttt{Flip Angle = 180$^\circ$}, \texttt{Averages = 2} &
High in-plane SNR, but slight smooth intensity shading across the field of view. \\
\textbf{Volume matrix} &
\texttt{256 $\times$ 256 $\times$ 128} voxels &
128 transaxial slices along the long axis of the fruit (\texttt{8,388,608} voxels total). \\
\textbf{Voxel spacing} &
\texttt{0.2344 $\times$ 0.2344 $\times$ 0.5000\,mm$^3$} ($\approx 0.0275\,\text{mm}^3/\text{voxel}$) &
\textbf{Anisotropic voxels:} slice thickness (\texttt{0.5\,mm}) is more than twice the in-plane pixel size (\texttt{0.234\,mm}), causing partial-volume averaging between slices. \\
\textbf{Files in \href{\ghtree/Lab4-Comp-Imag/Kiwi-challenge/data}{\texttt{data/}}} &
\href{\ghblob/Lab4-Comp-Imag/Kiwi-challenge/data/kiwi_20050217191129_131073_RARE_high_res_axial.nii.gz}{\texttt{kiwi\_20050217...high\_res\_axial.nii.gz}}\newline
\href{\ghblob/Lab4-Comp-Imag/Kiwi-challenge/data/kiwi_20050217191129_131073_RARE_high_res_axial.json}{\texttt{kiwi\_20050217...high\_res\_axial.json}}\newline
\href{\ghblob/Lab4-Comp-Imag/Kiwi-challenge/data/kiwi_micro_gt_labels.nii.gz}{\texttt{kiwi\_micro\_gt\_labels.nii.gz}} \& \href{\ghblob/Lab4-Comp-Imag/Kiwi-challenge/data/kiwi_micro_gt_roi.json}{\texttt{...\_roi.json}}\newline
\href{\ghtree/Lab4-Comp-Imag/Kiwi-challenge/data/dicom}{\texttt{dicom/}} (128 files) \& \href{\ghblob/Lab4-Comp-Imag/Kiwi-challenge/data/kiwi_064.png}{\texttt{kiwi\_064.png}} &
Ready for Python (\texttt{nibabel}, \texttt{pydicom}, \texttt{SimpleITK}, \texttt{MONAI}, \texttt{scikit-image}) as well as \href{https://www.slicer.org/}{3D Slicer}, \href{http://www.itksnap.org/}{ITK-SNAP}, or \href{https://imagej.net/software/fiji/}{Fiji/ImageJ}. Includes a reproducible semi-automated (\textit{``silver-standard''}) \texttt{80 $\times$ 80 $\times$ 16} \textbf{3D Micro-Ground-Truth} label volume (\texttt{98} 3D seed instances, generated by \href{\ghblob/Lab4-Comp-Imag/Kiwi-challenge/scripts/curate_micro_gt.py}{\texttt{scripts/curate\_micro\_gt.py}}) aligned to the scan affine. \\
\end{xltabular}
}

\subsection{Notebook Sequence \& Interactive Workshop}

The notebooks form a step-by-step progression and run in the main course \texttt{uv} environment (\texttt{bmed365-2027} kernel) or directly in Google Colab \textit{(in Colab, \texttt{resolve\_kiwi\_file(...)} and \texttt{resolve\_kiwi\_dicom\_dir()} download the DICOM, NIfTI, and Micro-GT files automatically from the \texttt{main} branch)}. Every notebook includes tiered \textbf{MSc \& PhD Learning Outcomes}, \textbf{Imaging $\leftrightarrow$ Clinical Glossaries}, \textbf{paired intuitive \& rigorous mathematical definitions}, \textbf{post-figure diagnostic interpretations}, and collapsible \textbf{``Check Your Understanding'' Q\&As}:

{\footnotesize
\renewcommand{\arraystretch}{1.28}
\rowcolors{2}{rowalt}{white}
\begin{xltabular}{\textwidth}{L{4.25cm} Y C{1.45cm}}
\toprule
\rowcolor{uibblue!12}\textbf{Notebook} & \textbf{Description} & \textbf{Launch} \\
\midrule
\endfirsthead
\toprule
\rowcolor{uibblue!12}\textbf{Notebook} & \textbf{Description} & \textbf{Launch} \\
\midrule
\endhead
\bottomrule
\endfoot
\href{\ghblob/Lab4-Comp-Imag/Kiwi-challenge/notebooks/00-kiwi-dicom-to-nifti.ipynb}{\texttt{00-kiwi-dicom-to-nifti.ipynb}} &
\textbf{Part 0:} History of medical image formats (\texttt{XITE/BIFF} $\to$ \texttt{DICOM} $\to$ \texttt{NIfTI} \& \texttt{BIDS}), reading the 128 Bruker 7\,T DICOM slices, Larmor $B_0$ \& privacy tags, homogeneous $4\times 4$ \texttt{RAS+} affine synthesis, avoiding 3 classic geometry bugs, \& converting \texttt{DICOM $\to$ NIfTI} (\texttt{pydicom}, \texttt{SimpleITK}, \texttt{dcm2niix}). &
\colablink{\colabbase/Lab4-Comp-Imag/Kiwi-challenge/notebooks/00-kiwi-dicom-to-nifti.ipynb} \\

\href{\ghblob/Lab4-Comp-Imag/Kiwi-challenge/notebooks/01-kiwi-explore.ipynb}{\texttt{01-kiwi-explore.ipynb}} &
\textbf{Part 1:} Volumetric NIfTI inspection, optical vs.\ 7\,T $T_2$-weighted RARE spin-echo contrast, anisotropic voxel geometry (\texttt{0.234 $\times$ 0.234 $\times$ 0.500\,mm$^3$}), partial-volume physics, \& $4\times 4$ \texttt{sform} affine world coordinates. &
\colablink{\colabbase/Lab4-Comp-Imag/Kiwi-challenge/notebooks/01-kiwi-explore.ipynb} \\

\href{\ghblob/Lab4-Comp-Imag/Kiwi-challenge/notebooks/02-kiwi-seeds-explore.ipynb}{\texttt{02-kiwi-seeds-explore.ipynb}} &
\textbf{Part 2:} Quantitative 16-bit vs.\ 8-bit LUT integrity, 2D stereological overcounting (\textit{Wicksell's fallacy}), Otsu vs.\ morphological attribute (\texttt{diameter\_closing}) top-hat filtering, standardized multispectral $K$-means, \& the 2011--2012 Dagstuhl / CUHK convex Mumford--Shah milestone. &
\colablink{\colabbase/Lab4-Comp-Imag/Kiwi-challenge/notebooks/02-kiwi-seeds-explore.ipynb} \\

\href{\ghblob/Lab4-Comp-Imag/Kiwi-challenge/notebooks/03-kiwi-paradigms-and-verification.ipynb}{\texttt{03-kiwi-paradigms-and-\newline verification.ipynb}} &
\textbf{Part 3 (Lab 4 Live Walkthrough):} Full 3D segmentation ladder (Rung 1 $\to$ Rung 2B $\to$ Foundation Models \& Agents), quantitative ``lesion load'', Curated 3D Micro-GT scoring (\texttt{P/R/F1} + calibrated count), \& the \textbf{4-Panel Semi-Verifiability Dashboard}. &
\colablink{\colabbase/Lab4-Comp-Imag/Kiwi-challenge/notebooks/03-kiwi-paradigms-and-verification.ipynb} \\

\href{\ghblob/Lab4-Comp-Imag/Kiwi-challenge/notebooks/03b-kiwi-interactive-gt-and-interrater.ipynb}{\texttt{03b-kiwi-interactive-gt-\newline and-interrater.ipynb}} &
\textbf{Part 3b (Interactive GT \& Inter-Rater Workshop):} In-notebook 3D Jupyter-Widget seed delineator (native anisotropic, \texttt{SimpleITK} isotropic resampled, or Silver-GT template), personal \texttt{MyStuff/} NIfTI export, \textbf{human inter-rater variability}, \textbf{Semantic vs.\ Instance ($\text{PQ} = \text{SQ} \times \text{RQ}$)}, \& \textbf{STAPLE consensus}. &
\colablink{\colabbase/Lab4-Comp-Imag/Kiwi-challenge/notebooks/03b-kiwi-interactive-gt-and-interrater.ipynb} \\

\href{\ghblob/Lab4-Comp-Imag/Kiwi-challenge/notebooks/04-kiwi-semantic-vs-instance.ipynb}{\texttt{04-kiwi-semantic-vs-\newline instance.ipynb}} &
\textbf{Part 4 (Semantic, Instance \& Panoptic Deep Dive):} Formal definitions \& label-permutation symmetry, head-to-head \textbf{3D PyTorch Semantic U-Net vs.\ Multi-Task Distance/Cellpose-Flow Instance U-Net}, \& fast offline prompt-and-track baseline. &
\colablink{\colabbase/Lab4-Comp-Imag/Kiwi-challenge/notebooks/04-kiwi-semantic-vs-instance.ipynb} \\

\href{\ghblob/Lab4-Comp-Imag/Kiwi-challenge/notebooks/05-kiwi-promptable-foundation-models.ipynb}{\texttt{05-kiwi-promptable-\newline foundation-models.ipynb}} &
\textbf{Part 5 (Promptable Vision Foundation Models Workshop):} Live \textbf{\texttt{SAM 2.1} vs.\ \texttt{MedSAM}} inference via Hugging Face \texttt{transformers} --- ViT token-stride physics \& \textbf{Patch-Zooming}, box-only vs.\ multi-mask ambiguity, \textbf{1-Click \texttt{Sam2VideoModel} 3D slice-memory propagation} (unbounded vs.\ physically gated), \& \textbf{zero-training neurosymbolic whole-kiwi seed census}. &
\colablink{\colabbase/Lab4-Comp-Imag/Kiwi-challenge/notebooks/05-kiwi-promptable-foundation-models.ipynb} \\

\href{\ghblob/Lab4-Comp-Imag/Kiwi-challenge/notebooks/06-kiwi-frontier-ai-paradigms.ipynb}{\texttt{06-kiwi-frontier-ai-\newline paradigms.ipynb}} &
\textbf{Part 6 (Frontier Biomedical AI Capstone):} Completing the 6-pillar SOTA taxonomy in pure PyTorch --- \textbf{3D Self-Supervised Learning (\texttt{Models Genesis} / \texttt{MAE} + 34-param few-shot probing)}, \textbf{Selective State-Space Models (\texttt{3D U-Mamba} via $\mathcal{O}(\log L)$ parallel associative scan)}, \textbf{Conditional Flow-Matching Diffusion (\texttt{MedSegDiff} / \texttt{CIMD} posterior count uncertainty $P(N \mid \mathbf{X})$)}, \& \textbf{Differentiable Neurosymbolic AI (\texttt{Logic Tensor Networks} + \texttt{Lab 5} hard-core split curing whole-organ \texttt{OOD} collapse)}. &
\colablink{\colabbase/Lab4-Comp-Imag/Kiwi-challenge/notebooks/06-kiwi-frontier-ai-paradigms.ipynb} \\

\href{\ghblob/Lab4-Comp-Imag/Kiwi-challenge/notebooks/07-kiwi-neurosymbolic-ai.ipynb}{\texttt{07-kiwi-neurosymbolic-\newline ai.ipynb}} &
\textbf{Part 7 (Differentiable Neurosymbolic AI Deep Dive):} Focused educational deep dive into \texttt{NB 06}'s winning paradigm --- empirically auditing botanical, 7\,T $T_2$ core/rim, \& \texttt{Lab 5} hard-core priors on Micro-GT, proving the \texttt{InstanceNorm3d} coordinate-erasure trap, deriving LTN $p$-mean quantifiers ($p=3$), running a controlled \textbf{5-model ablation} with an explicit \textbf{Bayes Log-Odds \texttt{Organ-GPS}} + leakage-free threshold calibration (\texttt{0} polar cap FPs, \texttt{-99.92\%} OOD voxels), auditing neural--symbolic division of labor \& error anatomy, and translating every lesson to clinical medical imaging (\texttt{CMB}, \texttt{MS}, \texttt{BraTS}, \texttt{IMC}). &
\colablink{\colabbase/Lab4-Comp-Imag/Kiwi-challenge/notebooks/07-kiwi-neurosymbolic-ai.ipynb} \\
\end{xltabular}
}

\begin{kiwibox}[Shared Lab 4 $\leftrightarrow$ Lab 5 Visual Anatomy, Histology \& MRI Contrast Atlas]
Wonder \textit{what} every bright and dark zone in the 7\,T scan is microscopically --- or \textit{why} the white central \textit{columella} turns \textbf{dark} in $T_2$ MRI while translucent \textit{locular mucilage} turns \textbf{brightest}? Open the companion illustrated atlas \href{\ghblob/Lab5-Comp-Mod/Kiwi-challenge/notebooks/00a-kiwi-anatomical-and-physiological-atlas.ipynb}{\textbf{\texttt{00a-kiwi-anatomical-and-physiological-atlas.ipynb}}} (\colablink{\colabbase/Lab5-Comp-Mod/Kiwi-challenge/notebooks/00a-kiwi-anatomical-and-physiological-atlas.ipynb}) in \href{\ghblob/Lab5-Comp-Mod/Kiwi-challenge/README.md}{\texttt{Lab5-Comp-Mod/Kiwi-challenge/}}.
\end{kiwibox}

\begin{infobox}[Methodological Deep Dives \& Technical Reports (\texttt{docs/})]
\begin{itemize}
    \item \href{\ghblob/Lab4-Comp-Imag/Kiwi-challenge/docs/neurosymbolic-report/README.md}{\textbf{\texttt{docs/neurosymbolic-report/README.md}}} (\href{\ghblob/Lab4-Comp-Imag/Kiwi-challenge/docs/neurosymbolic-report/neurosymbolic-kiwi-report.pdf}{\textbf{Master's-Level Technical Report PDF}} $|$ \href{\ghblob/Lab4-Comp-Imag/Kiwi-challenge/docs/neurosymbolic-report/neurosymbolic-kiwi-report.tex}{\textbf{\LaTeX}} $|$ \href{\ghblob/Lab4-Comp-Imag/Kiwi-challenge/docs/neurosymbolic-report/references.bib}{\textbf{BibTeX}}): Comprehensive IMRAD \textbf{``Thesis + Code'' technical report} expanding Notebook~\texttt{07} with 3-seed $\times$ 6-model ablations, causal \texttt{Organ-GPS} interpretability, mass-effect axiom stress tests (\texttt{0--8\,mm}), explicit silver-standard GT provenance \& circularity audit, 33 clickable peer-reviewed references, \textbf{Technical Appendices A--F} (including manual \texttt{freeview}/\texttt{ITK-SNAP} GT editing procedures), and \textbf{5 modular PyTorch companion notebooks} (\texttt{01}--\texttt{05}) backed by \texttt{kiwi\_ns\_lib.py}.
    \item \href{\ghblob/Lab4-Comp-Imag/Kiwi-challenge/docs/promptable-foundation-models.md}{\textbf{\texttt{docs/promptable-foundation-models.md}}} (\href{\ghblob/Lab4-Comp-Imag/Kiwi-challenge/docs/promptable-foundation-models.pdf}{\textbf{Printable PDF}}): Rigorous MSc/PhD guide and Apple Silicon \texttt{mps} benchmark of SOTA promptable vision foundation models (\texttt{SAM 2.1}, \texttt{MedSAM}, \texttt{MedSAM2}, \texttt{MicroSAM}, \texttt{nnInteractive}, \texttt{SAM 3}) --- starting from a non-technical visual introduction and LLM-to-Vision Transformer bridge, and covering ViT patch-stride physics, Patch-Zooming, 2D-to-3D slice memory drift, and ready-to-run \texttt{transformers} recipes.
    \item \href{\ghblob/Lab4-Comp-Imag/Kiwi-challenge/docs/frontier-ai-paradigms.md}{\textbf{\texttt{docs/frontier-ai-paradigms.md}}} (\href{\ghblob/Lab4-Comp-Imag/Kiwi-challenge/docs/frontier-ai-paradigms.pdf}{\textbf{Printable PDF}}): Companion MSc/PhD guide to Notebook~\texttt{06} opening with a non-technical visual introduction and covering \textbf{3D Self-Supervised Learning (\texttt{SSL})}, \textbf{Selective State-Space Models (\texttt{U-Mamba} / \texttt{SegMamba})}, \textbf{Conditional Diffusion \& Flow Matching (\texttt{MedSegDiff} / \texttt{CIMD} / \texttt{AnoDDPM})}, \textbf{Differentiable Neurosymbolic AI (\texttt{Logic Tensor Networks})}, plus \textbf{Implicit Neural Fields (\texttt{INRs})}, \textbf{Graph Neural Networks (\texttt{GNNs})}, and \textbf{Object-Query Decoders (\texttt{Mask2Former})}.
\end{itemize}
\end{infobox}

\begin{infobox}[Runtime \& Platform Notes]
\begin{itemize}
    \item \textbf{Notebooks \texttt{00}--\texttt{03b} (CPU-only):} Run on any CPU in \texttt{$\sim$5--50\,s} each with modest memory use (\texttt{$< 2$\,GB} RAM). No GPU is needed.
    \item \textbf{Notebook \texttt{04} (PyTorch 3D U-Net, zero external downloads):} Automatically picks \texttt{cuda} (NVIDIA GPU), \texttt{mps} (Apple Silicon GPU), or \texttt{cpu}. Training both 3D U-Nets (\texttt{$\sim$340k} parameters each, 95 epochs) takes \textbf{\texttt{$\sim$3\,s} on CUDA} (Ubuntu + NVIDIA RTX A5000) and \textbf{\texttt{$\sim$11\,s} on Apple Silicon MPS} (MacBook Pro M4 Max); CPU fallback also works if no GPU is available.
    \item \textbf{Notebook \texttt{05} (\texttt{SAM 2.1} \& \texttt{MedSAM} via \texttt{transformers}):} Downloads ungated checkpoints (\texttt{facebook/sam2.1-hiera-small}, $\sim$184\,MB, and \texttt{wanglab/medsam-vit-base}, $\sim$375\,MB) on first run and executes all 2D/3D ablations + whole-fruit scan in \textbf{\texttt{$\sim$45\,s} on Apple Silicon MPS} or CUDA / Colab T4 GPU.
    \item \textbf{Notebook \texttt{06} (Frontier AI Capstone --- 100\% pure PyTorch, zero external downloads):} Runs all four frontier paradigms (\texttt{3D SSL}, \texttt{3D U-Mamba}, \texttt{Conditional Flow-Diffusion}, \texttt{LTN Neurosymbolic}) + whole-fruit OOD audit in \textbf{\texttt{$\sim$70\,s} on Apple Silicon MPS} or CUDA / Colab T4 GPU, with zero extra packages beyond \texttt{pyproject.toml}.
    \item \textbf{Notebook \texttt{07} (Differentiable Neurosymbolic AI Deep Dive --- 100\% pure PyTorch, zero external downloads):} Runs empirical prior audits, \texttt{InstanceNorm3d} erasure proof, 5-model controlled ablation (\texttt{SSL} warm-start + \texttt{Baseline}, \texttt{CoordConv3D}, $p=1$, $p=3$, \& \texttt{Train-Calibrated Bayes Organ-GPS + Hard-Core Split}), and mechanistic interpretability audits in \textbf{\texttt{$\sim$95\,s} on Apple Silicon MPS} or CUDA / Colab T4 GPU.
    \item \textbf{Google Colab:} All nine notebooks download the required DICOM/NIfTI/Micro-GT files automatically from \texttt{main} via \texttt{resolve\_kiwi\_file(...)} / \texttt{resolve\_kiwi\_dicom\_dir()}. Default CPU runtime is fine for \texttt{00}--\texttt{03b}; for \texttt{04}--\texttt{07}, select \textbf{Runtime $\to$ Change runtime type $\to$ T4 GPU}.
\end{itemize}
\end{infobox}

% =====================================================================
\section{A Spectrum of Approaches Across the Course}
% =====================================================================

As we progress through BMED365, consider how each family of methods handles (or fails on) the kiwi seeds:

\begin{enumerate}
    \item \textbf{Classical Rule-Based \& Signal Processing (2D vs.\ 3D):} Global or adaptive thresholding, background bias-field correction, morphological opening/closing, distance-transform watershed, Laplacian/Difference of Gaussians (LoG/DoG) blob detection, or 3D ellipsoidal template matching. \textit{Watch out for:} Dark peel voxels, radial vascular bundles, and adjacent seeds merging into single large connected components.
    \item \textbf{Classical Machine Learning \& Radiomics (cf.\ \href{\ghblob/Lab0-ML/README.md}{Lab 0}):} Stack filtered volumes (intensity, gradient magnitude, Hessian/blobness at seed scale, distance from fruit center) into a multi-feature matrix for unsupervised clustering ($K$-means, GMM) or supervised pixel/voxel classification (Random Forest / XGBoost, or interactive tools like \href{https://www.ilastik.org/}{ilastik}).
    \item \textbf{Spatial Graphs \& Structural Priors (cf.\ \href{\ghblob/Lab1-NetworkSci-PSN/README.md}{Lab 1}):} Kiwifruit seeds are not scattered at random: they are arranged inside bright radial locules surrounding the central fibrous core (\textit{columella}, which looks white optically but appears dark in $T_2$-weighted MRI). Representing candidate seed centroids as a 3D spatial graph (\texttt{GNN} / \texttt{torch\_geometric}) helps spot outliers and merged clusters.
    \item \textbf{Supervised 3D ConvNets, Self-Supervised Learning \& Foundation Models (cf.\ \href{\ghtree/Lab2-DL}{Lab 2}, \href{\ghblob/Lab4-Comp-Imag/README.md}{Lab 4}, \href{\ghblob/Lab4-Comp-Imag/Kiwi-challenge/docs/promptable-foundation-models.md}{\texttt{promptable-foundation-models}}, \& \href{\ghblob/Lab4-Comp-Imag/Kiwi-challenge/docs/frontier-ai-paradigms.md}{\texttt{frontier-ai-paradigms}}):} Train a compact 2D/3D U-Net (\texttt{NB 04}), pretrain a \textbf{3D Masked Context Self-Supervised (\texttt{SSL})} encoder across unlabeled whole-organ blocks (\texttt{NB 06} Step~1), apply biomedical instance segmenters (\href{https://github.com/stardist/stardist}{StarDist}, \href{https://www.cellpose.org/}{Cellpose}), or test promptable vision foundation models (\href{https://github.com/facebookresearch/sam2}{SAM 2}, \href{https://github.com/bowang-lab/MedSAM}{MedSAM}, \href{https://github.com/computational-cell-analytics/micro-sam}{MicroSAM}, \href{https://huggingface.co/MIC-DKFZ/nnInteractive}{nnInteractive} in \texttt{NB 05}).
    \item \textbf{Selective State-Space Models (\texttt{U-Mamba}) \& Generative Conditional Diffusion (\texttt{NB 06} Steps~2--3):} Capture $\mathcal{O}(N)$ global volumetric context at fine CNN resolution via \textbf{Tri-Oriented 3D \texttt{U-Mamba} / \texttt{SegMamba}} ($\mathcal{O}(\log L)$ associative scan), and quantify \textbf{aleatoric boundary \& count uncertainty} $P(N \mid \mathbf{X})$ on partial-volume and kissing seeds via \textbf{Conditional Flow-Matching / Diffusion (\texttt{MedSegDiff} / \texttt{CIMD})}.
    \item \textbf{Differentiable Neurosymbolic \& Prior-Constrained AI (\texttt{NB 05} Step~4, \texttt{NB 06} Step~4, \& \texttt{NB 07} Deep Dive):} Upgrade from pipeline-level botanical masking to \textbf{Logic Tensor Networks (\texttt{LTN})}: compiling anatomical radial compartments, Bloch $T_2$ contrast invariants, and \texttt{Lab 5} Mat\'ern hard-core steric repulsion ($d_{\min} \ge 0.88\,\text{mm}$) into differentiable real-logic losses on unlabeled slices to eliminate spatial out-of-distribution (\texttt{OOD}) false positives.
    \item \textbf{AI-Assisted Coding \& Agentic Workflows (cf.\ \href{\ghblob/Lab-QuickStart/README.md}{Quick Start} \& \href{\ghtree/Lab3-GenAI-LLM}{Lab 3}):} Ask a multimodal LLM directly \textit{``How many seeds are in \texttt{kiwi\_064.png}?''} and observe why raw visual counting fails. Then use an \textbf{AI coding partner or agent} (Gemini, Claude, ChatGPT, Cursor) to write a 3D Python pipeline, inspect diagnostic overlay plots, check physical constraints, and iteratively refine the algorithm.
\end{enumerate}

% =====================================================================
\section{Evaluation \& Semi-Verifiability: How Do We Know an Answer Is Right?}
\label{sec:eval}
% =====================================================================

In many real biomedical problems --- and in AI-assisted coding --- running code that produces \textit{a number} is easy; \textbf{verifying} that the number reflects biological reality is the hard part. Because complete voxel-by-voxel ground truth does not exist for this scan, the Kiwi Challenge is \textbf{semi-verifiable} through multiple independent checks:

\begin{itemize}
    \item \textbf{Curated 3D Sub-Volume Audit (``Micro-Ground-Truth'' --- Silver $\to$ Gold Standard):} \href{\ghblob/Lab4-Comp-Imag/Kiwi-challenge/data/kiwi_micro_gt_labels.nii.gz}{\nolinkurl{data/kiwi_micro_gt_labels.nii.gz}} and \href{\ghblob/Lab4-Comp-Imag/Kiwi-challenge/data/kiwi_micro_gt_roi.json}{\nolinkurl{data/kiwi_micro_gt_roi.json}} provide a reproducible \textbf{semi-automated (``silver-standard'')} \texttt{80 $\times$ 80 $\times$ 16} equatorial validation sector ($x=[128,208)$, $y=[128,208)$, $z=[56,72)$, \texttt{98} 3D seed instances, \texttt{75} strictly interior) generated by \href{\ghblob/Lab4-Comp-Imag/Kiwi-challenge/scripts/curate_micro_gt.py}{\nolinkurl{scripts/curate_micro_gt.py}} (padded multi-scale 3D LoG + anisotropic grey closing + watershed + physical-mm $K$-means pair splitting, refined across multi-slice visual audits). Notebook~\texttt{03} computes 3D instance \textbf{Precision, Recall, F1}, and a \textbf{bias-calibrated whole-fruit seed count} ($\hat{N}_{\text{cal}} = \hat{N} \cdot \text{Precision} / \text{Recall}$). Because the \texttt{uint16} NIfTI shares the exact \texttt{256 $\times$ 256 $\times$ 128} grid and affine of the MRI scan, you can inspect or manually refine it into a human-certified \textbf{gold standard} in 4 quick steps:
    \begin{enumerate}
        \item \textbf{Load main MRI:} Open \href{\ghblob/Lab4-Comp-Imag/Kiwi-challenge/data/kiwi_20050217191129_131073_RARE_high_res_axial.nii.gz}{\nolinkurl{data/kiwi_20050217191129_131073_RARE_high_res_axial.nii.gz}} as the \textit{Main Image} in \href{http://www.itksnap.org/}{\textbf{ITK-SNAP}} (\textit{File $\to$ Open Main Image}) or \href{https://www.slicer.org/}{\textbf{3D Slicer}} (\textit{Add Data $\to$ Volume}).
        \item \textbf{Overlay 3D seed labels:} Load \href{\ghblob/Lab4-Comp-Imag/Kiwi-challenge/data/kiwi_micro_gt_labels.nii.gz}{\nolinkurl{data/kiwi_micro_gt_labels.nii.gz}} as a \textbf{Segmentation} (\textit{ITK-SNAP: Segmentation $\to$ Open Segmentation}; \textit{3D Slicer: Add Data $\to$ Description = Segmentation}).
        \item \textbf{Jump to the Micro-GT sector:} Navigate the cursor to voxel \textbf{\texttt{(x=168, y=168, z=64)}} (center of \texttt{x=128..207, y=128..207, z=56..71}) and click \textit{Update 3D} / \textit{Show 3D} to inspect all 98 labeled 3D seeds across axial, coronal, and sagittal views.
        \item \textbf{Verify, edit, or add a second wedge (in ITK-SNAP/Slicer, FreeSurfer \texttt{freeview}, or Notebook~\texttt{03b}):} Use the 3D paintbrush/polygon tool --- or launch FreeSurfer \texttt{freeview} directly on macOS via \href{\ghblob/Lab4-Comp-Imag/Kiwi-challenge/scripts/view_kiwi_freeview.sh}{\nolinkurl{scripts/view_kiwi_freeview.sh}} (\texttt{--with-pred --3d} also overlays the bit-aligned Neurosymbolic whole-fruit 3D prediction with custom \href{\ghblob/Lab4-Comp-Imag/Kiwi-challenge/scripts/kiwi_seeds_ColorLUT.txt}{\nolinkurl{kiwi_seeds_ColorLUT.txt}}, and \texttt{--gt <path>} opens a personal working copy), or use the zero-install Jupyter-Widget 3D tri-planar editor in \href{\ghblob/Lab4-Comp-Imag/Kiwi-challenge/notebooks/03b-kiwi-interactive-gt-and-interrater.ipynb}{\nolinkurl{03b-kiwi-interactive-gt-and-interrater.ipynb}} --- to refine boundary voxels, save personal rater NIfTI masks to \texttt{MyStuff/kiwi-gt/}, measure \textbf{inter-rater Dice / Cohen's $\kappa$ / Instance F1}, and build a multi-rater \textbf{STAPLE consensus} standard.
    \end{enumerate}
    \item \textbf{Physical Volume Sanity Checks:} At \texttt{0.0275\,mm$^3$/voxel}, a realistic kiwifruit seed ($\sim$1--3\,mm$^3$) spans roughly \textbf{35--110 voxels}. Plotting the histogram of connected-component volumes immediately reveals tiny noise specks ($<10$ voxels) and fused multi-seed blobs ($>200$ voxels).
    \item \textbf{Anatomical Zone Constraints:} Seeds only exist in the inner pericarp ring --- zero detections should fall in the central fibrous columella, the outer skin, or air outside the fruit.
    \item \textbf{3D Multi-Planar Consistency:} A real seed spans $\sim$3--5 consecutive axial slices and remains compact when viewed in \textbf{coronal} and \textbf{sagittal} cross-sections.
    \item \textbf{Perturbation \& Cross-Method Consensus:} Does changing a threshold by 5\% double your seed count? Do two very different methods (e.g., 3D LoG blob detection vs.\ a deep instance segmenter) agree on the same seed locations?
\end{itemize}

\subsection{Prospective Wet-Lab \& In-Situ Ground Truth: How Would You Measure True Seed Count, Mass, and Volume?}

Manual voxel painting in \texttt{ITK-SNAP}, \texttt{3D Slicer}, or \texttt{STAPLE} (\href{\ghblob/Lab4-Comp-Imag/Kiwi-challenge/notebooks/03b-kiwi-interactive-gt-and-interrater.ipynb}{\texttt{03b}}) only measures \textbf{human image interpretation (``silver standard'')} --- it is still bounded by the scanner's \texttt{0.50\,mm} slice thickness and partial-volume physics. In clinical imaging research, true \textbf{gold-standard validation} requires an \textbf{orthogonal physical, chemical, or histological measurement} (\textit{radiology--pathology correlation}).

If you scan a fresh kiwifruit tomorrow to benchmark a segmentation pipeline, distinguish \textbf{two complementary classes of physical ground truth}:
\begin{itemize}
    \item \textbf{Spatially Co-Registered 3D Ground Truth ($\{(x_i, y_i, z_i, v_i)\}$):} Preserves 3D seed coordinates inside the intact berry so every MRI prediction can be matched 1-to-1 (\texttt{Dice}, \texttt{Panoptic Quality}, \texttt{3D Instance F1}).
    \item \textbf{Global / Bulk Wet-Lab Ground Truth ($N_{\text{true}}, M_{\text{seeds}}, V_{\text{seeds}}, V_{\text{kiwi}}$):} Destructive \textit{ex-situ} extraction sacrifices $(x, y, z)$ positions, \textbf{but} delivers near-exact whole-organ seed count and parenchymal ``lesion load'' ($V_{\text{seeds}} / V_{\text{kiwi}}$) to expose systematic counting or partial-volume bias.
\end{itemize}

{\footnotesize
\renewcommand{\arraystretch}{1.26}
\rowcolors{2}{rowalt}{white}
\begin{xltabular}{\textwidth}{L{2.55cm} L{4.25cm} Y L{3.35cm}}
\toprule
\rowcolor{uibblue!12}\textbf{Regime \& Discipline} & \textbf{How to Measure $N_{\text{seeds}}$, $M_{\text{seeds}}$, or $V_{\text{seeds}}$} & \textbf{What It Captures \& Key Metrological Pitfalls to Avoid} & \textbf{Clinical / Biomedical Validation Parallel} \\
\midrule
\endfirsthead
\toprule
\rowcolor{uibblue!12}\textbf{Regime \& Discipline} & \textbf{How to Measure $N_{\text{seeds}}$, $M_{\text{seeds}}$, or $V_{\text{seeds}}$} & \textbf{What It Captures \& Key Metrological Pitfalls to Avoid} & \textbf{Clinical / Biomedical Validation Parallel} \\
\midrule
\endhead
\bottomrule
\endfoot
\textbf{In-Situ Physical}\newline \textit{(Non-destructive 3D)} &
\textbf{3D X-Ray Micro-CT (\texttt{10--25\,$\mu$m} isotropic)} \textit{(optionally contrast-enhanced via Lugol's $\text{I}_2/\text{KI}$ or \texttt{PTA})} &
Orthogonal electron-density contrast ($\rho_{\text{testa}} \approx 1.18\,\text{g/cm}^3$ vs.\ $\rho_{\text{gel}} \approx 1.03\,\text{g/cm}^3$) at $\sim$1,000$\times$ finer voxel volume than 7\,T MRI; rigid/affine \texttt{MRI $\leftrightarrow$ $\mu$CT} registration yields true 3D voxel/instance GT. \textit{(Optical/OCT fails past $\sim$2--3\,mm due to Mie scattering; hard seeds \& hairy peel air pockets shadow ultrasound.)} &
\textit{Ex vivo} specimen micro-CT or synchrotron \texttt{HiP-CT} validation of \textit{in vivo} clinical MRI. \\

\textbf{Semi-In-Situ Physical}\newline \textit{(Coordinate-preserving)} &
\textbf{Fiducial-Guided Cryo-Sectioning / 3D-Printed Slicing Jig + Block-Face Photography} &
Freeze berry with 3 MRI-visible fiducial tubes, slice at \texttt{0.5--1.0\,mm} parallel to the MRI axial plane, and photograph each exposed block face. \textbf{Trap:} \textit{Wicksell's corpuscle fallacy (1925)} --- a $\sim$2.2\,mm seed appears on 3--4 slices! Use \textbf{Sterio's physical disector} (count a seed \textit{only} on the slice where its top pole first appears). &
Patient-specific \textbf{3D-printed prostatectomy, mastectomy, or glioma slicing molds} co-registering \texttt{mpMRI} to whole-mount \texttt{H\&E}. \\

\textbf{Ex-Situ Biological / Enzymatic}\newline \textit{(Cleanest extraction!)} &
\textbf{Pectinase + Cellulase Digestion} (\textit{Aspergillus niger} food-grade enzymes, \texttt{pH 4.0--4.5}, \texttt{40\,$^\circ$C}, \texttt{45\,min}) &
Sticky locular \textbf{pectic mucilage} (\texttt{polygalacturonan}) glues seeds to pulp during mechanical sieving. Pectinase + cellulase liquefy the hydrogel and thin parenchymal walls into clear fluid while leaving \textbf{heavily lignified sclereid seed coats (\textit{testa}) 100\% intact}. \textit{(24--48\,h natural fermentation also works, but risks early water imbibition.)} &
\textbf{Collagenase / dispase tissue dissociation} for single-cell FACS counting; \textbf{acid (\texttt{HCl}) maceration} for whole-kidney glomerular counting. \\

\textbf{Ex-Situ Chemical}\newline \textit{(Clearing \& viability sorting)} &
\textbf{Warm \texttt{1--2\%} $\text{Na}_2\text{CO}_3$ or \texttt{EDTA} Wash + \texttt{500\,$\mu$m} Mesh Sieve + Sucrose Flotation ($\rho \approx 1.10\,\text{g/cm}^3$)} &
Chelates $\text{Ca}^{2+}$ ``egg-box'' cross-links in pectin so pulp flushes cleanly through a \texttt{500\,$\mu$m} sieve (mature seeds $\sim$1.2 $\times$ 2.3\,mm are 100\% retained). In $\rho \approx 1.10\,\text{g/cm}^3$ sucrose/glycerol, \textbf{viable filled seeds ($\rho \approx 1.18\,\text{g/cm}^3$) sink} while hollow/aborted seed coats float. &
\texttt{Ficoll} / \texttt{Percoll} density-gradient cell separation \& chemical tissue decellularization. \\

\textbf{Ex-Situ Physical Metrology}\newline \textit{(Exact $N, M, V$)} &
\textbf{1.\ Count $N_{\text{true}}$:} Backlit \texttt{600\,dpi} flatbed scan of dispersed non-touching monolayer (+ 100-seed gravimetric check $N \approx 100\,M_{\text{tot}}/m_{100}$)\newline
\textbf{2.\ Mass $M_{\text{seeds}}$:} Analytical balance ($\pm 0.1\,\text{mg}$)\newline
\textbf{3.\ Volume $V_{\text{seeds}}$:} \textbf{Ethanol / heptane liquid pycnometry} + 2D optical profilometry ($v_i = \frac{4}{3}\pi a_i b_i^2$) &
\textbf{Three classic wet-lab metrology traps:}\newline
$\bullet$ \textbf{Hydration state ($M$ \& $V$):} Record \textbf{blotted ``saturated surface-dry'' (\texttt{SSD}) mass} ($\sim$1.5--1.9\,mg/seed) immediately after brief extraction to match \textit{in-situ} MRI hydration, before oven-drying (\texttt{60\,$^\circ$C, 24\,h}, $\sim$1.1--1.4\,mg/seed).\newline
$\bullet$ \textbf{Why ethanol ($\rho = 0.789\,\text{g/cm}^3$) or heptane instead of water?} Water's high surface tension traps micro-air bubbles inside the honeycomb pits of the reticulated \textit{testa} (inflating $V_{\text{seeds}}$) \textit{and} triggers osmotic swelling.\newline
$\bullet$ \textbf{Total vs.\ individual volume:} Pycnometry gives exact bulk $V_{\text{seeds}}$, while \texttt{600\,dpi} silhouette major/minor axes give the full single-seed distribution $v_i$ (\texttt{median} \& \texttt{IQR}). &
Archimedes whole-organ volumetry, gas/liquid pycnometry, \& high-throughput digital pathology morphometry. \\
\end{xltabular}
}

\begin{protocolbox}[Hands-On Block 2 Student Benchmark Protocol (``Scan $\to$ Segment $\to$ Pectinase Digest $\to$ Monolayer Scan \& Ethanol Pycnometry'', $< 90$\,min wet lab)]
\begin{enumerate}
    \item \textbf{Whole-organ physical baseline (\textit{in situ}):} Weigh the intact kiwifruit ($\pm 0.01\,\text{g}$) and measure true whole-fruit volume $V_{\text{kiwi}}$ by Archimedes water displacement (briefly submerging the intact berry in a graduated cylinder / overflow vessel) \textbf{before} scanning --- directly validating your 3D MRI foreground mask.
    \item \textbf{Acquire 3D image \& lock predictions:} Scan the intact berry (preclinical 7\,T MRI, clinical 3\,T MRI, or fiducial-guided serial optical slices) and pre-register your algorithm's predicted $\hat{N}$, $\hat{V}_{\text{kiwi}}$, $\hat{V}_{\text{seeds}}$, and lesion load $\hat{V}_{\text{seeds}}/\hat{V}_{\text{kiwi}}$ \textit{before} opening the fruit (\textit{blinding against experimenter bias}).
    \item \textbf{Enzymatic release:} Halve the berry, scoop all pulp into \texttt{100\,mL} warm water (\texttt{40\,$^\circ$C}, \texttt{pH $\sim$4.2}) with \texttt{0.5--1\%} food-grade wine/juice \textbf{pectinase} (\texttt{45\,min} gentle stirring), and rinse over a \textbf{\texttt{500\,$\mu$m} sieve} until 100\% clean seeds remain.
    \item \textbf{SSD mass \& ethanol pycnometry ($M_{\text{seeds}}, V_{\text{seeds}}$):} Gently blot seeds to \textbf{saturated surface-dry (\texttt{SSD})} state on lint-free paper, weigh total $M_{\text{seeds,SSD}}$ ($\pm 0.1\,\text{mg}$), and measure displaced volume $V_{\text{seeds}}$ in a calibrated \texttt{10\,mL} \textbf{ethanol pycnometer} at \texttt{20\,$^\circ$C}:
    \[
        V_{\text{seeds}} = \frac{m_{\text{pyc+EtOH}} + M_{\text{seeds,SSD}} - m_{\text{pyc+EtOH+seeds}}}{\rho_{\text{EtOH}}(20\,^\circ\text{C})}
    \]
    \item \textbf{Monolayer optical census ($N_{\text{true}}$ \& $\{v_i\}$):} Spread all seeds without touching on a backlit glass/acrylic plate, scan at \texttt{600\,dpi}, count exact $N_{\text{true}}$ ($\pm 1$ seed) + prolate-spheroid volumes $v_i = \frac{4}{3}\pi a_i b_i^2$, and audit where your 3D MRI pipeline over- or under-estimated reality.
\end{enumerate}
\end{protocolbox}

% =====================================================================
\section{How We Use This Challenge in BMED365}
% =====================================================================

\begin{itemize}
    \item \textbf{Lecture 1 $\to$ Lab 4 Thought Experiment:} Keep this dataset in mind throughout Block~1 as you learn ML, networks, deep learning, and prompt/agentic engineering. In \textbf{Lab 4 (Mon 8 Feb 2027)}, we compare approaches side-by-side --- from naive thresholding failures to modern foundation-model and agentic pipelines --- and examine how each method is evaluated.
    \item \textbf{Lab 5 Companion Module --- \href{\ghblob/Lab5-Comp-Mod/Kiwi-challenge/README.md}{Computational Morphogenesis \& 7T MRI Digital Twin (\nolinkurl{Lab5-Comp-Mod/Kiwi-challenge})}:} In Lab~5, we reverse the direction (\texttt{Biology \& Physics $\to$ 3D Form \& Synthetic MRI}): modeling seed-hormone allometry, \texttt{XY} floral sex determination, 160-day turgor/carbohydrate ODEs, 1D Ring Turing locule patterning (\texttt{K = 33}), 2D paracrine diffusion PDEs, 3D spatial point-process repulsion (\texttt{R = 1.19}), Tikhonov PDE inverse problems, and a \textbf{generative 3D Bloch-MRI Digital Twin} with known ground-truth seeds --- accompanied by two printable non-technical + methodological guides (\href{\ghblob/Lab5-Comp-Mod/Kiwi-challenge/docs/foundations-atlas-and-morphogenesis.pdf}{Part 1 PDF} \& \href{\ghblob/Lab5-Comp-Mod/Kiwi-challenge/docs/inverse-problems-and-digital-twin.pdf}{Part 2 PDF}).
    \item \textbf{Optional \href{\ghblob/Individual-project/README.md}{Individual Project} Topic (Block 2):} If you are looking for a concrete, self-contained dataset for your Block~2 digital poster and speed-poster presentation, you are welcome to use the Kiwi Challenge --- either from the \textbf{3D segmentation \& validation} side (\texttt{Lab 4}) or the \textbf{mechanistic morphogenesis \& Digital Twin} side (\href{\ghblob/Lab5-Comp-Mod/Kiwi-challenge/README.md}{\texttt{Lab 5}}).
\end{itemize}

% =====================================================================
\section{Additional Background \& Visualizations}
% =====================================================================

\begin{itemize}
    \item \textbf{From 1990s \texttt{XITE/BIFF} (\texttt{<filename>.us}) to \texttt{DICOM}, \texttt{NIfTI}, and \texttt{BIDS}:} \href{\ghblob/Lab4-Comp-Imag/Kiwi-challenge/notebooks/00-kiwi-dicom-to-nifti.ipynb}{\texttt{00-kiwi-dicom-to-nifti.ipynb}} traces how biomedical image formats evolved from early custom Unix/C pipelines --- such as unpacking 12-bit/pixel frames from the 1992 \textbf{1.0\,T Siemens Magnetom Impact} scanner at Haukeland University Hospital into 16-bit unsigned short \textbf{BIFF} files (\texttt{<filename>.us}) via UiO's \href{https://www.mn.uio.no/ifi/english/research/groups/dsb/resources/software/xite/}{\textbf{XITE}} library (\texttt{xshow} + C / MATLAB 4.0; \href{https://www.tandfonline.com/doi/abs/10.3109/00207459509004883}{Lundervold et al., \textit{Int.\ J.\ Neurosci.}, 1995}; \href{https://www.mn.uio.no/ifi/english/research/groups/dsb/resources/software/xite/usersmanual.pdf}{XITE User's Manual}) --- to clinical \textbf{DICOM 3.0} (\texttt{LPS}) and research \textbf{NIfTI-1} (\texttt{RAS+}) + \textbf{BIDS} sidecars.
    \item \textbf{Dagstuhl / CUHK 3D Variational Segmentation (2011--2012):} Following a November 2011 Dagstuhl Seminar discussion with Prof.\ Raymond H.\ Chan (The Chinese University of Hong Kong), PhD student Xiaohao Cai applied their two-stage convex Mumford--Shah variational segmentation method (cf.\ Cai, Chan \& Zeng, \textit{SIAM J.\ Imaging Sci.}, 2013) to this 3D kiwi dataset --- see \href{\ghblob/Lab4-Comp-Imag/Kiwi-challenge/assets/kiwi_3d_cai_chan_2012.jpg}{\texttt{./assets/kiwi\_3d\_cai\_chan\_2012.jpg}} and the historical discussion in \href{\ghblob/Lab4-Comp-Imag/Kiwi-challenge/notebooks/02-kiwi-seeds-explore.ipynb}{\texttt{02-kiwi-seeds-explore.ipynb}}.
    \item \textbf{3D Volume Rendering (Dept.\ of Informatics, UiB):} Example 3D rendering from the UiB visualization group in \href{\ghblob/Lab4-Comp-Imag/Kiwi-challenge/assets/kiwi-stian-eikeland-vis-video-to-gif.png}{\texttt{./assets/kiwi-stian-eikeland-vis-video-to-gif.png}} (visualizations originally created by Ove Daae Lampe, Stian Eikeland, and Veronika \v{S}olt\'eszov\'a).
    \item \textbf{Botanical, Histological \& Physiological Reference Atlas:} \href{\ghblob/Lab5-Comp-Mod/Kiwi-challenge/notebooks/00a-kiwi-anatomical-and-physiological-atlas.ipynb}{\textbf{\nolinkurl{00a-kiwi-anatomical-and-physiological-atlas.ipynb}}} (macroscopic optical vs.\ 3-plane 7\,T MRI contrast inversion, $\mu\text{m}$-scale histology from \href{https://pmc.ncbi.nlm.nih.gov/articles/PMC3797463/}{Guo et al., \textit{Front.\ Plant Sci.}\ 2013 (\texttt{PMC3797463})}, dioecious \texttt{XY} life cycle, \& \texttt{BBCH} ontogeny from \href{https://pmc.ncbi.nlm.nih.gov/articles/PMC3261216/}{Richardson et al., \textit{BMC Plant Biol.}\ 2011 (\texttt{PMC3261216})}) $\;\vert\;$ \href{https://www.zespri.com/en-IN/blogdetail/all-about-kiwi-seeds}{All about kiwi seeds (Zespri)} $\;\vert\;$ \href{https://en.wikipedia.org/wiki/Kiwifruit}{Kiwifruit (\textit{Actinidia deliciosa}) morphology}.
\end{itemize}

\end{document}
